Un alambre tiene la forma de la hélice
\[
\vec r(t)=(a\cos t,a\sin t,bt),\qquad 0\leq t\leq2\pi,
\]
donde $a,b>0$ son longitudes. Su densidad es $\rho(x,y,z)=\rho_0(1+z/b)$, con $\rho_0>0$ en unidades de masa por longitud. Determine su masa y compárela con la de un alambre de la misma forma y densidad constante $\rho_0$.
